\section{Referencial teórico}
\label{sec:referencial}

A exposição ocupacional à inteligência artificial permite organizar a análise das diferenças entre ocupações, mas não deve ser interpretada automaticamente como substituição de trabalhadores. \citeonline{felten2021} apresentam uma medida de exposição que relaciona aplicações de inteligência artificial e capacidades utilizadas nas ocupações. Para esta pesquisa, essa abordagem oferece uma referência para comparar grupos ocupacionais, preservando a distinção entre exposição potencial, adoção efetiva e resultados observados no emprego.

Outra dimensão relevante é a possibilidade de realizar atividades a distância. \citeonline{dingel2020} examinam quais ocupações podem ser desempenhadas em casa, fornecendo uma referência para considerar diferenças na organização do trabalho. A discussão sobre automação e novas tarefas também ajuda a formular hipóteses sobre mudanças na demanda por trabalho \parencite{acemoglu2019}. Essas contribuições orientam o desenho da pesquisa, sem determinar antecipadamente o sinal ou a magnitude das associações que serão examinadas no caso brasileiro.

Este capítulo contém texto provisório para o exercício ia-6.1. A revisão completa deverá aprofundar os conceitos, examinar os limites das medidas e explicitar a relação entre cada referência e as hipóteses da dissertação.
\subsection{Medida de exposição ocupacional}

O trecho a seguir apresenta a medida de exposição ocupacional proposta pelos autores. A transcrição foi mantida no idioma original:

\begin{citacao}
We create and validate a new measure of an occupation's exposure to AI that we call the AI Occupational Exposure (AIOE). We use the AIOE to construct a measure of AI exposure at the industry level, which we call the AI Industry Exposure (AIIE) and a measure of AI exposure at the county level, which we call the AI Geographic Exposure (AIGE). We also describe several ways in which the AIOE can be used to create firm level measures of AI exposure.
\parencite[p.~2195]{felten2021}
\end{citacao}

Nesta pesquisa, a medida de exposição serve como referência para diferenciar ocupações. Sua utilização exige explicitar a correspondência entre classificações ocupacionais e preservar a distinção entre exposição potencial e adoção efetiva da tecnologia.